\documentclass{beamer}
\usepackage{amsmath}
\usepackage[version=4]{mhchem}
\usepackage{aurora}

% Your own mechanism shorthands: \newcommand works inside scenes.
% \attack{nucleophile atom}{target atom}: lone pair attacks an atom.
% \leave{atom}{leaving atom}: that bond's electrons go to the leaving atom.
\newcommand{\attack}[2]{\CurlyArrow[from=lp:#1, to=#2]}
\newcommand{\leave}[2]{\CurlyArrow[from=#1-#2, to=#2]}

\title{Reaction Mechanisms}
\author{Your Name}
\date{\today}

\begin{document}

\begin{frame}
  \titlepage
\end{frame}

\section{SN2}

% Atom-map numbers ([C:2]) name atoms: curly arrows and transforms refer to them.
\begin{frame}{Backside attack}
\begin{scene}[size={16,5.5}]
  \Molecule[id=nu, at={-5,0}, lp=auto]{[OH-:1]}
  \Molecule[id=sub, at={0,0}, lp={4}, with]{[CH3:2][Br:4]}
  \pause
  \attack{1}{2}
  \leave{2}{4}
  \pause
  \Transform[id={nu,sub}, lp={1:2,4:4}]{[OH:1][CH3:2].[Br-:4]}
  \ReactionArrow[from={-1.2,-2}, to={1.2,-2}, above={\ce{H2O}}]
\end{scene}
\end{frame}

\begin{frame}{Equations in chemistry notation}
  \[ \ce{CH3Br + OH- -> CH3OH + Br-} \]
  \pause
  Rate law: $\text{rate} = k\,[\ce{CH3Br}][\ce{OH-}]$
\end{frame}

\end{document}
